\section{Introduction}
\label{sec:introduction}

Fix one nonzero integer shift \(h_0\).  The high-row branch of the
fixed-shift program reaches a finite determinant coefficient
\begin{equation}
 a_X=(A_{C_X}(n))_{n\in\cN_X},
 \qquad C_X=\lfloor J_X\rfloor,
 \label{eq:intro-coefficient}
\end{equation}
formed after literal aggregation of a complete set of native
prime--M\"obius keys.  Its zero mode and determinant energy are
\begin{equation}
 Z_X:=\sum_{n\in\cN_X}A_{C_X}(n),
 \qquad
 D_X:=\sum_{n\in\cN_X}|A_{C_X}(n)|^2.
 \label{eq:intro-ZD}
\end{equation}
The zero-frequency concentration is
\begin{equation}
 \mathfrak F_{h_0}(X)
 :=
 \begin{cases}
 |Z_X|^2/D_X,&D_X>0,\\
 0,&D_X=0.
 \end{cases}
 \label{eq:intro-flatness}
\end{equation}
The zero convention in the second line is harmless only because then
\(a_X=0\) and hence \(Z_X=0\).
Here ``zero frequency'' means the normalized determinant-DFT
frequency \(r=0\) used in TPC-32.  It is neither a generic nonzero
determinant frequency nor any orbit, divisor, or transfer-operator
zero mode.

TPC-32 proved a conditional transfer at
\begin{equation}
 Q_X=X^{267/400+o(1)},
 \qquad
 J_X=X^{133/400+o(1)}:
 \label{eq:intro-scales}
\end{equation}
the estimate
\begin{equation}
 \mathfrak F_{h_0}(X)\le X^{1/400+o(1)}
 \label{eq:intro-threshold}
\end{equation}
delivers every strict saving below \(133/400\) for its complete
matched raw cutoff shell \citep{WangTPC32}.  TPC-33 isolated the
sufficient amplitude estimate
\begin{equation}
 |Z_X|\ll X^{o(1)}Q_X^2,
 \label{eq:intro-amplitude}
\end{equation}
but explicitly did not prove a matching lower bound for \(D_X\)
\citep{WangTPC33}.  Indeed, \eqref{eq:intro-amplitude} alone says
nothing about \eqref{eq:intro-threshold} when the denominator may be
small.

The purpose of this paper is to identify exactly what a completion
certificate must additionally prove before it controls \(D_X\).  The
finite completion, sparse-incidence, and scalar Perron layers that
motivate this interface were developed in
\citep{WangTPC75,WangTPC77,WangTPC80}; none of those layers by itself
identifies the determinant coefficient.
The
relevant normalized statistic is
\begin{equation}
 \boxed{\displaystyle
 \mathfrak D_{h_0}(X)
 :=\frac{D_X}{Q_X^3J_X^2}.}
 \label{eq:intro-energy-statistic}
\end{equation}
At the scales \eqref{eq:intro-scales}, the denominator
\(Q_X^3J_X^2\) is the energy scale that pairs with
\eqref{eq:intro-amplitude}.  The needed zero-loss statement is
\begin{equation}
 \mathfrak D_{h_0}(X)\ge X^{-o(1)}.
 \label{eq:intro-energy-target}
\end{equation}
No such physical estimate is proved here.

\subsection{The three interfaces}

There are three logically different maps:
\[
 \text{native coefficients}
 \xrightarrow{\ \cA_X+\cM_X\ }
 \text{pre-bin atoms}
 \xrightarrow{\ \cJ_X\ }
 \text{determinant bins}
 \xrightarrow{\ \bm 1^*\ }
 \text{zero mode}.
\]
The first arrow is a completion/reassembly problem.  The second can
lose energy through collisions in a determinant fiber.  The third is
the arithmetic signed sum.  Success at one arrow does not imply
success at either later arrow.

For the first arrow we prove a missing-Gram majorant theorem.  If
\(\cB_X=\cA_X+\cM_X\), then for every coefficient vector \(z\) with
\(\cA_Xz\ne0\),
\begin{equation}
 \left(1-\sqrt{\Xi_X(z)}\right)_+^2
 \le
 \frac{\|\cB_Xz\|_2^2}{\|\cA_Xz\|_2^2}
 \le
 \left(1+\sqrt{\Xi_X(z)}\right)^2,
 \label{eq:intro-specific}
\end{equation}
where
\begin{equation}
 \Xi_X(z)
 :=
 \widehat\kappa_X
 \frac{\langle D_{\cM,X}z,z\rangle}
      {\|\cA_Xz\|_2^2},
 \qquad
 \widehat\kappa_X=1+\rho(P_X).
 \label{eq:intro-Xi}
\end{equation}
Here \(D_{\cM,X}\) is the exact diagonal of the missing Gram and
\(P_X\) is an entrywise nonnegative majorant of its normalized
off-diagonal part.

For the second arrow we require a literal binning map, not an
analogy.  If \(s_X\) is a verified lower modulus of \(\cJ_X\) on the
actual reassembled subspace and
\[
 r_X=a_X-\cJ_X\cB_Xz_X,
\]
then the lower bound displayed in the abstract follows.  The residual
\(r_X\) records every mismatch between the completion carrier and the
physical determinant coefficient.

For the final arrow we prove the exact exponent arithmetic.  If
\begin{equation}
 |Z_X|\le X^{-\eta_Z+o(1)}Q_X^2,
 \qquad
 D_X\ge X^{-\lambda_D+o(1)}Q_X^3J_X^2,
 \label{eq:intro-exponents}
\end{equation}
then
\begin{equation}
 \mathfrak F_{h_0}(X)
 \le
 X^{1/400+\lambda_D-2\eta_Z+o(1)}.
 \label{eq:intro-exponent-transfer}
\end{equation}
Consequently these recorded bounds recover the TPC-32 threshold
whenever \(2\eta_Z\ge\lambda_D\), and this condition is sharp for the
two bounds alone.  The current amplitude target has \(\eta_Z=0\);
within this transfer ledger it therefore permits no positive
certified energy loss.  Such a loss cannot be absorbed unless it is
offset by a new amplitude saving.

\subsection{Level ledger}

The level labels used throughout are:
\begin{itemize}[leftmargin=2em]
\item \textbf{L0:} exact finite-dimensional algebra, including Gram
majorization, reassembly inequalities, and determinant-binning
identities;
\item \textbf{L1:} a verified literal interface from those identities
to the actual fixed-\(h_0\) coefficient, together with asymptotic
bounds for every certificate statistic;
\item \textbf{L2:} a new arithmetic estimate for the actual
fixed-\(h_0\) packet, and then a separately verified global block
cover.
\end{itemize}
This paper proves L0 theorems and states auditable L1 implications.
It proves no new L2 cancellation estimate.  In particular, none of
the following can be inferred from the results below: a fixed-shift
prime-pair asymptotic, positivity at \(h_0=2\), twin primes, or a
breach of the sieve parity barrier \citep{HeathBrown1982Parity}.

\subsection{Organization}

\Cref{sec:literal-packet} defines the physical packet and its global
normalization.  \Cref{sec:completion-survival} proves the
coefficient-specific Gram bound, and \cref{sec:uniform-reassembly}
extracts a uniform survival certificate.
\Cref{sec:dictionary} adds the determinant dictionary.
\Cref{sec:zero-slack} gives the exponent theorem.
\Cref{sec:nonidentifiability} supplies sharp no-go models.
\Cref{sec:full-block} states the conditional global transfer, and
\cref{sec:audit} gives the archive and stop protocol.
